% Revize 2026-10-27: cílené doplnění; nové úlohy a řešení jsou vlastní tvorba.
% Metodický vzor: 4_rocnik/pripravy/Uvod_do_kombinatoriky.pdf (samostatný dokument).
% Hodina 27, 1. ročník SPŠ. Rozsah: Rozpis_vyuky_MAT_1_rocnik.tex, řádek 27.

% Plán: 1_rocnik/tematicke_plany/SPS.pdf; společná projektová šablona.

% Učebnice: ucebnice/01_Zakladni_poznatky/Mocniny_a_odmocniny/Mocniny a Odmocniny.tex

% Původ úloh: Vlastní zadání a vlastní formulace výkladu; metodická inspirace zdroji.

% Metodický zdroj: externi_zdroje/realisticky/Mocniny/03_Vzorce_pro_mocniny_II.pdf

% URL: http://www.realisticky.cz/ucebnice/01%20Matematika%20S%C5%A0/01%20Z%C3%A1kladn%C3%AD%20poznatky/06%20Mocniny/03%20Vzorce%20pro%20mocniny%20II.pdf

% Metodický zdroj: externi_zdroje/realisticky/Mocniny/04_KISS.pdf

% URL: http://www.realisticky.cz/ucebnice/01%20Matematika%20S%C5%A0/01%20Z%C3%A1kladn%C3%AD%20poznatky/06%20Mocniny/04%20KISS.pdf

\documentclass[11pt,a4paper]{article}

% Spolecny styl pro ucitelske pripravy MAT (Work/MAT)
% Verze 1.1, 2026-10-23
\NeedsTeXFormat{LaTeX2e}
\ProvidesPackage{priprava}[2026/10/23 v1.1 Sablona ucitelskych priprav MAT]

% Zakladni typografie a cestina
\RequirePackage[T1]{fontenc}
\RequirePackage[utf8]{inputenc}
\RequirePackage[czech]{babel}
\RequirePackage{lmodern}
\RequirePackage{microtype}
\RequirePackage{amsmath,amssymb}
\RequirePackage{xcolor}
\RequirePackage{enumitem}
\RequirePackage[most]{tcolorbox}
\RequirePackage{tabularx,booktabs,array}
\RequirePackage{fancyhdr}
\RequirePackage{needspace}
\RequirePackage{geometry}
\RequirePackage[hidelinks]{hyperref}

\geometry{a4paper,top=17mm,bottom=18mm,left=23mm,right=23mm}
\definecolor{notegray}{gray}{0.39}
\setlength{\parindent}{0pt}
\setlength{\parskip}{2.5pt plus 1pt minus 0.5pt}
\setlist[itemize]{topsep=2pt,itemsep=1pt,parsep=0pt,leftmargin=1.45em}
\setlist[enumerate]{topsep=2pt,itemsep=2pt,parsep=0pt,leftmargin=1.65em}
\setcounter{secnumdepth}{0}
\renewcommand{\arraystretch}{1.12}

% Zahlavi: vlevo rocnik | globalni tema, vpravo cislo: nazev hodiny
% Zapati zustava zcela prazdne.
\pagestyle{fancy}
\fancyhf{}
\renewcommand{\headrulewidth}{0.35pt}
\renewcommand{\footrulewidth}{0pt}
\setlength{\headheight}{14pt}
\newcommand{\nastavPripravu}[4]{%
  \fancyhead[L]{\small #1\enspace\textbar\enspace #2}%
  \fancyhead[R]{\small #3: #4}%
}

% Minimalisticke cernobile boxy. Kazda definice zvlast.
\tcbset{commonlesson/.style={enhanced,breakable=false,colback=white,colframe=black,
  colbacktitle=white,coltitle=black,boxrule=0.45pt,
  arc=0.8mm,outer arc=0.8mm,left=3mm,right=3mm,top=1.6mm,bottom=1.6mm,
  toptitle=1mm,bottomtitle=0.7mm,fonttitle=\bfseries,
  attach boxed title to top left={xshift=2.5mm,yshift=-1.15mm},
  boxed title style={colback=white,colframe=white,boxrule=0pt,left=1mm,right=1mm}}}
\newtcolorbox{zadaniBox}[1][]{commonlesson,title={Zadání},#1}
\newtcolorbox{prikladBox}[1][]{commonlesson,title={Příklad},#1}
\newtcolorbox{definiceBox}[1][]{commonlesson,title={Definice},#1}

% Spolecny box pro zadani a jeho reseni, bez mezery mezi castmi.
% Pouziti:
%   \begin{prikladReseniBox}
%     Zadani ulohy.
%     \tcblower
%     \begin{align*} ... \end{align*}
%   \end{prikladReseniBox}
% Spodni cast ma jemny sedy podklad a delici linku; lze v ni sazet
% vice radku matematiky. Obe casti maji jen jeden spolecny ramecek.
\tcbset{solutionlesson/.style={
  commonlesson,
  enhanced jigsaw,
  breakable,
  colbacklower=black!3,
  segmentation style={solid,draw=black!25,line width=0.35pt},
  fontlower=\small\color{notegray},
  before lower={\noindent\textbf{Řešení:}\enspace},
  leftlower=3mm,rightlower=3mm
}}
\newtcolorbox{prikladReseniBox}[1][]{solutionlesson,title={Příklad},#1}
\newtcolorbox{zadaniReseniBox}[1][]{solutionlesson,title={Zadání},#1}

% Pomucky pro autora pripravy: sedive poznamky; \reseni pro starsi pripravy
\newcommand{\poznamka}[1]{\par\begingroup\small\color{notegray}\itshape
  \textnormal{\textbf{Poznámka pro učitele.}}\enspace #1\par\endgroup}
\newcommand{\reseni}[1]{\par\begingroup\small\color{notegray}
  \textbf{Očekávané řešení:}\enspace #1\par\endgroup}
\newcommand{\blok}[1]{\needspace{5\baselineskip}\subsection*{#1}}
\newcommand{\vyrok}[1]{\uv{#1}}


\nastavPripravu{1. ročník}{Mocniny}{27}{Mocnina mocniny a součinu}

\begin{document}

\blok{Výklad a souvislosti}

Zápis $(a^3)^2$ obsahuje dva činitele $a^3$, tedy šest činitelů $a$. Proti tomu součin $a^3a^2$ obsahuje pět činitelů $a$. Rozlišuj vnější operaci dříve, než začneš upravovat exponenty.

\begin{definiceBox}[title={Připomenutí definice},breakable=false]
\textbf{Mocnina součinu} $(ab)^n$ je pro $n\in\mathbb N$ součin $n$ stejných činitelů $ab$.
\end{definiceBox}


Z opakovaného násobení a přeskupení činitelů dostáváme
\[
(a^m)^n=a^{mn},\qquad (ab)^n=a^nb^n,\qquad
\left(\frac ab\right)^n=\frac{a^n}{b^n}\quad(b\ne0).
\]
Zde jsou $m,n\in\mathbb N$. Na součet se poslední dvě pravidla nevztahují: například $(2+3)^2=25$, ale $2^2+3^2=13$.

\blok{Řešené příklady}

\begin{prikladReseniBox}[title={Příklad 1},breakable=false]
Uprav $(x^3)^4$ a $(2a^2b)^3$.
\tcblower
$(x^3)^4=x^{12}$; $(2a^2b)^3=2^3a^6b^3=8a^6b^3$. Umocňujeme i číselný koeficient.
\end{prikladReseniBox}


\begin{prikladReseniBox}[title={Příklad 2},breakable=false]
Uprav $\left(\frac{-3x^2}{2y}\right)^2$.
\tcblower
Pro $y\ne0$ je výsledek $\frac{9x^4}{4y^2}$. Záporné znaménko je součástí základu a při sudé mocnině zmizí.
\end{prikladReseniBox}


\blok{Samostatné procvičení}

\begin{zadaniReseniBox}[title={Procvičení},breakable=false]
\begin{enumerate}[label=\alph*)]
\item Uprav $(b^2)^5$.
\item Uprav $(-2xy^3)^3$.
\item Uprav $\left(\frac{2a^3}{b^2}\right)^2$.
\end{enumerate}
\tcblower
\begin{enumerate}[label=\alph*)]
\item $b^{10}$, protože exponenty násobíme.
\item $-8x^3y^9$; lichý exponent zachovává záporné znaménko.
\item $\frac{4a^6}{b^4}$ při $b\ne0$.
\end{enumerate}
\end{zadaniReseniBox}


\poznamka{Po každém vzorci nech žáky vyřešit krátký samostatný příklad. Chybný vzorec pro mocninu součtu nyní jen vyvrať číselným příkladem; správné dvojčlenné vzorce patří až do 49. hodiny.}


\blok{Navazující procvičování ve třídě}

\begin{zadaniReseniBox}[title={Základ: který vzorec použít jako první?},breakable=false]
\begin{enumerate}[label=\alph*)]
\item Uprav $(a^2)^3$, $(-a^2)^3$, $(-a)^6$. Vysvětli rozdíl zápisů.
\item Uprav $(3x^2y)^2$, $(-2ab^2)^4$ a $\left(\frac{2x^3}{3y^2}\right)^3$.
\item Uprav $\frac{(x^2)^3x^4}{x^7}$ a $\frac{(2a^2b)^2}{4ab^2}$.
\end{enumerate}
\tcblower
\begin{enumerate}[label=\alph*)]
\item $a^6$, $-a^6$, $a^6$. Ve druhém zápisu umocňujeme záporný součin $-a^2$ lichým exponentem.
\item $9x^4y^2$, $16a^4b^8$, $\frac{8x^9}{27y^6}$ s $y\ne0$ u podílu.
\item $x^3$ při $x\ne0$; $a^3$ při $a,b\ne0$. Podmínky určuje původní jmenovatel.
\end{enumerate}
\end{zadaniReseniBox}

\begin{zadaniReseniBox}[title={Rezerva: najdi neplatný krok},breakable=false]
\begin{enumerate}[label=\alph*)]
\item Žák píše $(2x^3)^2=2x^6$. Najdi chybu a oprav ji.
\item Porovnej $(2+3)^3$ a $2^3+3^3$. Lze pravidlo pro součin použít na součet?
\end{enumerate}
\tcblower
\begin{enumerate}[label=\alph*)]
\item Je třeba umocnit oba činitele: $2^2x^6=4x^6$.
\item $125\ne35$. Pravidlo platí pro součin, nikoli součet.
\end{enumerate}
\end{zadaniReseniBox}

\blok{Volitelná rezerva}

\begin{prikladReseniBox}[title={Rozšíření procvičení},breakable=false]
Uprav $\frac{(2^3\cdot3^2)^3}{(2\cdot3)^5}$.
\tcblower
$\frac{2^9 3^6}{2^5 3^5}=2^4\cdot3=48$. Krátit před vyčíslováním je úspornější.
\end{prikladReseniBox}



\poznamka{Základ: původní odvození a první nová skupina od jednoho pravidla ke složeným úpravám. Při kontrole požaduj pojmenovat vnější operaci a základ mocniny. Druhá skupina je rezerva s chybou k diskusi. Vzorce pro mocniny součtu zde neodvozujeme, jen ukážeme neplatnost záměny se součinem. Domácí čtveřici zadej po samostatném pokusu ve třídě.}

\blok{Domácí procvičení}

\begin{zadaniReseniBox}[title={Domácí úlohy},breakable=false]
\begin{enumerate}[label=\alph*)]
\item Uprav $(b^4)^3$ a $(-3a^2b)^3$.
\item Uprav $\left(\frac{4x^2}{y^3}\right)^2$.
\item Uprav $\frac{(3a^2)^2a^3}{9a^5}$.
\item Uprav $\frac{(2xy^2)^3}{4x^2y^4}$.
\end{enumerate}
\tcblower
\begin{enumerate}[label=\alph*)]
\item $b^{12}$ a $-27a^6b^3$.
\item $\frac{16x^4}{y^6}$, $y\ne0$.
\item $a^2$, $a\ne0$.
\item $2xy^2$, $x,y\ne0$.
\end{enumerate}
\end{zadaniReseniBox}

\end{document}
